% =====================================================================
%  DATOS DEL LIBRO  —  edita SOLO este fichero para personalizar la portada
%  (puedes generarlo automáticamente en la web: "Textos listos para copiar")
% =====================================================================

% Nombre temático: el MISMO que pondrás en el formulario de solicitud.
% El título del libro y del depósito en BURJC Digital será
% "Materiales de la asignatura <nombre temático>".
\newcommand{\nombreTematico}{Nombre de la Asignatura}

% Subtítulo libre (qué contiene el libro)
\newcommand{\subtitulo}{Guía de estudio, apuntes, presentaciones, ejercicios y exámenes}

% Autoría: nombres completos, tal y como aparecerán en BURJC Digital
\newcommand{\autores}{Nombre Apellido Apellido y Nombre Apellido Apellido}
% Lo mismo, pero un nombre por línea (para la portada)
\newcommand{\autoresPortada}{Nombre Apellido Apellido\\ Nombre Apellido Apellido}

% Grado(s) / máster(es) y nombre de la asignatura en cada uno
\newcommand{\titulaciones}{%
  Grado en XXXXX (Campus) --- Nombre de la asignatura\\
  Grado en YYYYY (Campus) --- Nombre de la asignatura%
}

\newcommand{\curso}{2026-2027}
\newcommand{\anio}{2026}

% Licencia: "by" (Atribución) o "by-sa" (Atribución-CompartirIgual)
\newcommand{\licenciaCodigo}{by-sa}

% Enlace al depósito. Hasta que BURJC Digital asigne el handle, déjalo así:
% la propia biblioteca lo incluye al procesar el depósito.
\newcommand{\handle}{https://burjcdigital.urjc.es}

% Serie de TV URJC con los vídeos/audios (déjalo vacío {} si no presentas)
\newcommand{\serieTV}{}
